% 高轨卫星抵近观测 A 题论文主文件
% 编译：在本目录运行 ./build.sh（先把 ../figures/*.svg 转成 PDF，再用 xelatex 编译两遍）
% Windows 上用 xelatex 编译时 ctex 会自动使用宋体、黑体；Linux 上自动使用 Fandol 字体。
\documentclass[a4paper,12pt]{ctexart}

\usepackage[a4paper,top=2.5cm,bottom=2.5cm,left=2.5cm,right=2.5cm]{geometry}
\usepackage{amsmath,amssymb,bm}
\usepackage{graphicx}
\usepackage{caption,subcaption}
\usepackage{booktabs}
\usepackage[hidelinks]{hyperref}

\graphicspath{{figures/}}
\IfFileExists{xurl.sty}{\usepackage{xurl}}{}
\pagestyle{plain}
\urlstyle{same}

% ---------------- 格式要求 ----------------
% 正文：宋体小四
\AtBeginDocument{\zihao{-4}}
\linespread{1.3}
% 一级标题：居中、三号、黑体；二级标题：左顶格、四号、黑体
\ctexset{
  section = {
    format = \centering\heiti\zihao{3},
    name = {,},
    number = \arabic{section},
    aftername = \quad,
    beforeskip = 18pt, afterskip = 12pt
  },
  subsection = {
    format = \raggedright\heiti\zihao{4},
    indent = 0pt,
    aftername = \quad,
    beforeskip = 12pt, afterskip = 6pt
  }
}
% 图题、表题：黑体五号
\DeclareCaptionFont{heiwu}{\heiti\zihao{5}}
\captionsetup{font=heiwu, labelsep=quad}
\captionsetup[sub]{font=heiwu, labelformat=parens}
% 公式：居中，按全文顺序编号
\numberwithin{equation}{section}

\newcommand{\aGEO}{a_{\mathrm{GEO}}}

\begin{document}

% ---------------- 题目 ----------------
\begin{center}
  {\heiti\zihao{3} 高轨卫星抵近观测任务的轨道机动规划模型}
\end{center}

% ---------------- 摘要 ----------------
\begin{center}
  {\heiti\zihao{4} 摘\quad 要}
\end{center}

本文围绕地球同步轨道（GEO）卫星对 30 颗高轨目标的抵近观测机动规划问题，从轨道相对运动几何与光学观测约束的物理机理出发，建立了统一的“时变可观测管道”抽象模型与两阶段字典序优化框架（先最大化观测目标数，再最小化累计速度增量），并提出了零数据先验的三维相对运动物理状态动态规划与全耦合序列凸优化求解体系。

\textbf{针对问题一（48~h 单星观测 B1）}，将光照角、侦测角双圆锥约束与 $10\sim 50$~km 球壳约束联合刻画为时变可观测管道，提出“候选观测时刻确定—截面瞄准点选取—脉冲初值估计与精修”三步法。利用轨道面穿越且地日夹角 $\beta \le 90^\circ$ 条件锁定 2 个背阳侧候选时刻，沿角平分线选取瞄准点，经 SLSQP 精修得到最优两脉冲方案：累计速度增量 \textbf{5.2128~m/s}（间隔 13.56~h），0.05~s 网格误差不超过 1.60~mm，经 ATK 4.2.0 独立复核全部满足约束。

\textbf{针对问题二（30 天单星 R1 巡游 30 颗目标）}，针对静态图整数规划（M2-C，松弛上界约 21 颗）忽略状态耦合、基线大邻域搜索（M2-B，176.31~m/s）耗时数小时的局限，提出零先验主创新模型 \textbf{M2-D}：离散层基于 Hill/Clohessy--Wiltshire 能量积分推导两脉冲三维解析转移代价并预计算查表，结合多起点全范围确定性变邻域下降（3D-CW-VND-SE-DP）求解物理状态动态规划，自动识别出 B19 与 B20 的同次零脉冲伴随飞越；连续层建立虚拟控制时刻—矢量联合序列凸优化（Joint $(\Delta\bm v, \Delta t)$ VC-SCP），在二阶锥规划中同步优化脉冲矢量与施加时刻。M2-D 在 24 核并行下仅用 \textbf{7.4~min} 即实现 \textbf{30/30 全覆盖}，累计速度增量降至 \textbf{159.33~m/s}（有效脉冲 44/58 次，单次最大 24.38~m/s，最小间隔 7078.9~s），较基线逐段方案（210.39~m/s）下降 24.3\%，较数小时大邻域搜索再降 16.98~m/s（9.6\%）。

\textbf{针对问题三（30 天 R1 与新星 R2 双星协同巡游）}，按相对经度将目标划分为西侧 10 颗（R1）与东侧 20 颗（R2）消除折返，取 R2 初始半长轴偏离上限 $\delta a = 5.90$~km，并将 M2-D 推广为 \textbf{M3-D（脉冲时刻—矢量—初始轨道全耦合序列凸优化）}。最终求得全覆盖最优方案：R1 消耗 \textbf{13.50~m/s}（较 M3-C 下降 25.6\%），R2 消耗 \textbf{14.61~m/s}（$a - \aGEO = -5.895$~km，$e = 0.000482 < 0.001$），两星合计 \textbf{28.11~m/s}，较 M3-C（33.46~m/s）再降 16.0\%，较问题二单星方案下降 \textbf{82.4\%}。

\textbf{模型检验与灵敏度分析表明}：三问全部 61 次观测在 0.05~s 加密网格（最坏距离误差 $\le 2.03$~mm）与 ATK 4.2.0 独立数值积分中 100\% 通过复核；问题三起主导作用的约束为半长轴偏离上限 $\delta a$ 而非偏心率上限 $e_{\max}$，方案稳健性良好。

\vspace{4pt}
\noindent{\heiti 关键词：}高轨卫星抵近观测；Clohessy--Wiltshire 相对运动；物理状态动态规划；序列凸优化；双星协同巡游


\clearpage

\section{问题重述}

\subsection{问题背景}

地球同步轨道带距地面约 35786~km，通信、导航增强、气象和预警等高价值卫星大多部署于此，这一区域的空间态势感知因而受到持续关注。地基光学和雷达设备受距离和大气条件限制，难以获得高轨目标的细节信息。由观测卫星机动到目标附近、在合适的距离和光照条件下近距离成像，是获取高轨目标精细特征的主要手段，这种方式称为抵近观测。美国太空军的 GSSAP 卫星就运行在同步带附近，具备交会与近距离操作能力，用于跟踪同步带上的人造目标并识别其特性\cite{gssap}。

抵近观测对相对几何有较严格的要求。观测卫星与目标的距离既不能太远，否则成像分辨率不足，也不能太近，否则存在碰撞风险。光照条件和观测方向同样受限，观测卫星需要处在目标被太阳照亮的一侧，并以地球为背景进行观测，这样才能获得清晰的图像。上述条件还必须连续满足一段时间，才能完成一次有效观测。张海涛等针对单个地球静止轨道目标研究了绕飞持续观测的轨迹规划，使成像卫星在整个绕飞过程中保持有利的光照和观测角\cite{zhang2024}。

当观测对象从单个目标扩展到多个目标时，问题还要决定访问的顺序和时间。Alfriend 等研究了同步轨道卫星在轨服务的最优规划\cite{alfriend2006}，Zhou 等研究了多颗同步轨道卫星加注任务的规划优化\cite{zhou2015}，王峥嵘等提出了同步轨道多目标飞掠成像的任务规划方法\cite{wang2025}。这类问题把离散的排序决策和连续的轨道机动耦合在一起，求解难度随目标数量迅速增加。另一方面，高轨卫星携带的推进剂有限，每一次轨道机动都会消耗卫星的在轨寿命。因此，如何以尽量少的速度增量依次抵近多颗目标并完成观测，是高轨抵近观测任务规划的核心问题。

\subsection{问题提出}

题目给出红方观测卫星 R1 和蓝方 30 颗目标卫星 B1 至 B30 在初始历元 2027 年 1 月 1 日 00:00:00 UTC 的位置和速度，坐标系为地心惯性系 ICRF。目标卫星均位于地球同步轨道带附近，全程不机动。轨道运动采用地球中心引力二体模型，引力常数 $\mu = 398600.4418~\mathrm{km^3/s^2}$，不考虑 $J_2$ 项、日月引力和光压等摄动，太阳只用于确定光照方向。

一次有效观测需要在至少 60~s 内连续满足三个条件。第一个条件是观测卫星与目标的距离保持在 10~km 至 50~km 之间。第二个条件是光照角不超过 $45^\circ$，光照角是以目标为顶点、指向太阳的射线与指向观测卫星的射线之间的夹角。第三个条件是侦测角不超过 $45^\circ$，侦测角是以目标为顶点、指向地心的射线与指向观测卫星的射线之间的夹角。两个角度都按有向射线计算，取值范围为 $0^\circ$ 至 $180^\circ$。观测卫星采用瞬时脉冲机动，每完成一次有效观测之前最多机动两次，单次速度增量不超过 30~m/s，同一颗卫星相邻两次机动的间隔不少于 3600~s，每颗观测卫星的累计速度增量不超过 300~m/s。在上述条件下，题目提出了以下三个问题。

{\heiti 问题一}\quad 要求在初始历元后 48~h 内，通过 R1 的轨道机动完成对目标 B1 的一次有效观测。方案需要给出每次机动的时刻和三维速度增量、各段自由飞行时间以及观测的起止时刻，并验证全部观测约束，使累计速度增量最小。

{\heiti 问题二}\quad 将任务扩展到 30 天，要求 R1 依次机动，对 30 颗目标尽可能多地完成有效观测。优化目标分为两级，首先使成功观测的不同目标数量最多，在此基础上再使 R1 的累计速度增量最小。

{\heiti 问题三}\quad 为弥补单颗卫星观测能力的不足，红方提前发射了一颗偏心率小于 0.001 的高轨卫星 R2，与 R1 一同在 30 天内观测这 30 颗目标。R2 在 $t_0$ 时刻的轨道由我方设计。R2 的机动约束和速度增量预算与 R1 相同，两颗卫星的预算相互独立。优化目标仍按两级排序，首先使两颗卫星合计观测的不同目标数量最多，再使两颗卫星的累计速度增量之和最小。

\section{问题分析}

三个问题共享相同的二体轨道动力学与单次光学观测几何条件，但决策结构逐层递进：从单目标两脉冲机动，扩展到单星多目标巡游排序，再进一步耦合双星任务分配与新星初始轨道设计。本节先分析单次观测的几何结构，再分析同步带附近轨道机动的代价特点，最后据此给出三问的建模思路。

\subsection{观测几何分析}

观测条件可以用目标周围的空间区域来理解。距离约束要求观测卫星位于以目标为球心、半径分别为 10~km 和 50~km 的两个球面之间。光照角约束要求观测卫星位于以目标指向太阳的射线为轴、半顶角为 $45^\circ$ 的圆锥内。侦测角约束要求观测卫星位于以目标指向地心的射线为轴、半顶角为 $45^\circ$ 的圆锥内。三者同时满足时观测卫星才处在可观测位置，观测卫星 R、目标卫星 B 与太阳、地心之间的几何关系如图~\ref{fig:geometry}(a) 所示。两个圆锥的轴随太阳和目标的运动不断转动，可观测区域因而是时变的。距离下界相当于在目标周围挖去了一个半径 10~km 的小球，区域中间是空的，不是凸集。这类区域上的优化问题没法直接按凸优化求解，第 6 节构造凸子问题时会用一个平面把这个空洞替换掉。

两个圆锥有公共部分的充要条件是两条轴线的夹角不超过 $90^\circ$。记目标处太阳方向与地心方向的夹角为 $\beta$，则可观测的必要条件为
\begin{equation}
  \beta(t) = \angle\big(\bm r_\odot(t) - \bm r_b(t),\ -\bm r_b(t)\big) \le 90^\circ .
  \label{eq:beta}
\end{equation}
太阳距离远大于同步轨道半径，这个条件近似等价于目标位于地球的背阳一侧。同步轨道目标每天绕地球一周，$\beta$ 也以一天为周期变化，每天大约有一半时间满足这个条件。图~\ref{fig:geometry}(b) 给出了目标 B1 在前 48~h 内 $\beta$ 的变化，阴影部分为满足必要条件的时段。这个条件只与目标和太阳有关，可以在求解之前预先算出每颗目标的可观测时段。

\begin{figure}[!htb]
  \centering
  \begin{subfigure}[b]{0.41\textwidth}
    \centering
    \includegraphics[width=\linewidth]{fig1a_geometry.pdf}
    \caption{观测几何关系}
  \end{subfigure}\hfill
  \begin{subfigure}[b]{0.57\textwidth}
    \centering
    \includegraphics[width=\linewidth]{fig1b_beta_b1.pdf}
    \caption{B1 处 $\beta$ 在前 48~h 内的变化}
  \end{subfigure}
  \caption{观测几何与可能观测时段}
  \label{fig:geometry}
\end{figure}

轨道面的差异进一步缩小了观测机会。R1 的轨道倾角为 $4.38^\circ$，30 颗目标的倾角分布在 $0.12^\circ$ 到 $9.01^\circ$ 之间。两颗卫星的轨道面相差 $1^\circ$ 时，法向相对运动的振幅约为 736~km，远大于 50~km 的观测距离上限。而在同步轨道上把轨道面转动 $1^\circ$ 约需 53.7~m/s 的速度增量，用 300~m/s 的预算逐一对准 30 颗目标的轨道面是不经济的。因此在低燃料的巡游方案中，观测卫星应优先利用两条轨道交线附近的自然近距离交会来完成观测。轨道交线每天经过两次，再叠加 $\beta \le 90^\circ$ 的条件，对于轨道面差异明显的目标，通常每天存在约一次低代价的观测机会。这一规律用来生成候选观测时刻，从而减少需要搜索的时段；求解时速度增量的方向不受限制，实际观测时刻由精确的轨道优化确定。

\subsection{机动代价分析}

所有卫星都在同步带附近，彼此之间的相对运动以沿轨道方向的相位漂移为主。设观测卫星的半长轴比同步轨道大 $\Delta a$，它相对同步轨道的经度漂移率为 $\dot\lambda = -\tfrac32 n\,\Delta a/a$，而在近圆轨道上改变半长轴 $\Delta a$ 需要的切向速度增量为 $\Delta v = \tfrac{v}{2a}\Delta a$。两式合起来得到
\begin{equation}
  \Delta v = \frac{v}{3n}\,\big|\Delta\dot\lambda\big| \approx 2.84~\mathrm{m/s}\ \text{每}\ 1^\circ/\text{天} .
  \label{eq:drift}
\end{equation}
要在 $D$ 天内追上相位相差 $\Delta\lambda$ 的目标，需要先建立 $\Delta\lambda/D$ 的漂移率，到达后再把它消除，代价约为 $5.68\,\Delta\lambda/D$~m/s。可见机动代价主要由相位差和可用时间决定。时间越充裕，所需的漂移率越小，代价也越低。

\begin{figure}[!htb]
  \centering
  \includegraphics[width=0.8\textwidth]{fig2_targets.pdf}
  \caption{初始时刻 30 颗目标在同步带上的分布（径向放大显示）}
  \label{fig:targets}
\end{figure}

图~\ref{fig:targets} 画出了初始时刻 30 颗目标在相位和半长轴上的分布。目标大致对称地分布在 R1 的东西两侧，各自延伸约 $19^\circ$。它们的半长轴都非常接近同步轨道，彼此漂移得很慢，一个月下来相对位置几乎不变。所以可以把这 30 颗目标看成一条基本固定的队列，R1 要做的是沿着这条队列依次走过去。由此可以得到两点判断。第一，访问顺序应沿相位单向推进，来回折返会让漂移的路程和燃料都成倍增加。第二，相位接近、轨道面也接近的目标可以在同一次飞越中先后观测，比如 B19 和 B20 几乎挨在一起，观测完一颗顺路就能观测另一颗，省去一次单独的转移。


\subsection{建模思路}

基于上述分析，三个问题可以沿“单次观测单元、单星序列巡游、双星选轨与分工”这条主线逐层建模，总体框架如图~\ref{fig:framework} 所示。


问题一是基本单元。给定 R1 的初始状态和目标 B1，决策变量是两次脉冲的时刻、两个速度增量矢量和观测开始时刻，约束是连续 60~s 内每一时刻都满足观测条件，这是一个带半无限约束的非线性规划。48~h 内 B1 有两个候选观测时段，越晚的时段留给相位调整的时间越多，所需速度增量越小。由于可观测区域非凸，求解时需要从多个初值出发搜索。约束要求 60~s 内每一时刻都满足，优化时每隔 1~s 检查一次，得到方案后再每隔 0.05~s 复查一遍，确认相邻检查时刻之间也没有越界。

\begin{figure}[!htb]
  \centering
  \includegraphics[width=0.78\textwidth]{fig3_framework.pdf}
  \caption{总体框架}
  \label{fig:framework}
\end{figure}

问题二把基本单元串成序列。决策同时包括离散部分和连续部分。离散部分是选哪些目标、按什么顺序、在哪一天的机会观测；连续部分是每次脉冲的时刻和矢量。前后两段转移共享中间的轨道状态，一段的机动会改变下一段的起点，所以逐段独立优化得到的方案并不是整体最优。对此，本文在离散层基于三维 Clohessy--Wiltshire（CW）能量积分两脉冲解析代价构建物理状态动态规划，并在全排列空间做确定性变邻域搜索以确定最优访问序列与交会日期；在连续层则将 30 天计划的全部脉冲矢量与施加时刻一并纳入序列凸优化做全计划联合精修，同时以传统束搜索和大邻域搜索作为基线对照。

问题三在问题二的外面增加了 R2 的轨道设计和目标分工。题面所说的 R2 提前发射，意味着它在 $t_0$ 时刻已经在轨，偏心率小于 0.001 约束的是 $t_0$ 时刻的初始轨道，这个初始轨道由我们设计。$t_0$ 之后的每次机动都计入 R2 的速度增量，机动后的轨道不再受偏心率约束。题目没有给出 R2 半长轴的具体范围，只说明它是高轨卫星。我们引入参数 $\delta a$，要求 R2 的初始半长轴与同步轨道半长轴之差不超过 $\delta a$，用来刻画“位于同步带内”这一要求。目标在相位上分布于 R1 的两侧，由两颗卫星各负责一侧并单向扫描，每颗卫星需要跨越的相位范围约减少一半。$\delta a$ 的取值和分工方案在第 7 节讨论。


% ---------------- 参考文献 ----------------
\begingroup
\renewcommand{\section}[2]{\begin{center}{\heiti\zihao{4}参考文献}\end{center}}
\zihao{-4}\raggedright
% 参考文献：按正文首次引用顺序编号，GB/T 7714 著录。[3]、[4] 的卷期页码待核对原文。
\begin{thebibliography}{9}
\setlength{\itemsep}{2pt}
\bibitem{gssap} United States Space Force. Geosynchronous Space Situational Awareness Program[EB/OL]. \url{https://www.spoc.spaceforce.mil/About-Us/Fact-Sheets/Display/Article/2381739/geosynchronous-space-situational-awareness-program}.
\bibitem{zhang2024} 张海涛, 王伟林, 张雅声, 等. 地球静止轨道航天器绕飞持续观测任务轨迹规划与控制[J]. 国防科技大学学报, 2024, 46(1): 74-86.
\bibitem{alfriend2006} ALFRIEND K T, LEE D J, CREAMER N G. Optimal servicing of geosynchronous satellites[J]. Journal of Guidance, Control, and Dynamics, 2006, 29(1): 203-206.
\bibitem{zhou2015} ZHOU Y, YAN Y, HUANG X, et al. Mission planning optimization for multiple geosynchronous satellites refueling[J]. Advances in Space Research, 2015, 56(11): 2612-2625.
\bibitem{wang2025} 王峥嵘, 曹晓瑞, 黄喜元, 等. 地球同步轨道多目标飞掠成像服务任务规划方法[J]. 航天控制, 2025, 43(3): 15-23.
\bibitem{kraft1988} KRAFT D. A software package for sequential quadratic programming[R]. Köln: DFVLR, 1988. DFVLR-FB 88-28.
\bibitem{scipy2020} VIRTANEN P, GOMMERS R, OLIPHANT T E, et al. SciPy 1.0: fundamental algorithms for scientific computing in Python[J]. Nature Methods, 2020, 17(3): 261-272.
\end{thebibliography}

\endgroup

\end{document}
